%% LyX 2.4.4 created this file.  For more info, see https://www.lyx.org/.
%% Do not edit unless you really know what you are doing.
\documentclass[english]{beamer}
\usepackage[T1]{fontenc}
\usepackage[latin9]{inputenc}
\setcounter{secnumdepth}{3}
\setcounter{tocdepth}{3}

\makeatletter
%%%%%%%%%%%%%%%%%%%%%%%%%%%%%% Textclass specific LaTeX commands.
% this default might be overridden by plain title style
\newcommand\makebeamertitle{\frame{\maketitle}}%
% (ERT) argument for the TOC
\AtBeginDocument{%
  \let\origtableofcontents=\tableofcontents
  \def\tableofcontents{\@ifnextchar[{\origtableofcontents}{\gobbletableofcontents}}
  \def\gobbletableofcontents#1{\origtableofcontents}
}

%%%%%%%%%%%%%%%%%%%%%%%%%%%%%% User specified LaTeX commands.
\beamertemplateshadingbackground{red!5}{structure!5}

\usepackage{beamerthemeshadow}
\usepackage{pgfnodes,pgfarrows,pgfheaps}
\usepackage{caption}

\beamertemplatetransparentcovereddynamicmedium

\pgfdeclareimage[width=0.6cm]{icsi-logo}{beamer-icsi-logo}
\logo{\pgfuseimage{icsi-logo}}

\newcommand{\Class}[1]{\operatorname{\mathchoice
  {\text{\small #1}}
  {\text{\small #1}}
  {\text{#1}}
  {\text{#1}}}}

\newcommand{\Lang}[1]{\operatorname{\text{\textsc{#1}}}}

% This gets defined by beamerbasecolor.sty, but only at the beginning of
% the document
\colorlet{averagebackgroundcolor}{normal text.bg}

\newcommand{\tape}[3]{%
  \color{structure!30!averagebackgroundcolor}
  \pgfmoveto{\pgfxy(-0.5,0)}
  \pgflineto{\pgfxy(-0.6,0.1)}
  \pgflineto{\pgfxy(-0.4,0.2)}
  \pgflineto{\pgfxy(-0.6,0.3)}
  \pgflineto{\pgfxy(-0.4,0.4)}
  \pgflineto{\pgfxy(-0.5,0.5)}
  \pgflineto{\pgfxy(4,0.5)}
  \pgflineto{\pgfxy(4.1,0.4)}
  \pgflineto{\pgfxy(3.9,0.3)}
  \pgflineto{\pgfxy(4.1,0.2)}
  \pgflineto{\pgfxy(3.9,0.1)}
  \pgflineto{\pgfxy(4,0)}
  \pgfclosepath
  \pgffill

  \color{structure}  
  \pgfputat{\pgfxy(0,0.7)}{\pgfbox[left,base]{#1}}
  \pgfputat{\pgfxy(0,-0.1)}{\pgfbox[left,top]{#2}}

  \color{black}
  \pgfputat{\pgfxy(-.1,0.25)}{\pgfbox[left,center]{\texttt{#3}}}%
}

\newcommand{\shorttape}[3]{%
  \color{structure!30!averagebackgroundcolor}
  \pgfmoveto{\pgfxy(-0.5,0)}
  \pgflineto{\pgfxy(-0.6,0.1)}
  \pgflineto{\pgfxy(-0.4,0.2)}
  \pgflineto{\pgfxy(-0.6,0.3)}
  \pgflineto{\pgfxy(-0.4,0.4)}
  \pgflineto{\pgfxy(-0.5,0.5)}
  \pgflineto{\pgfxy(1,0.5)}
  \pgflineto{\pgfxy(1.1,0.4)}
  \pgflineto{\pgfxy(0.9,0.3)}
  \pgflineto{\pgfxy(1.1,0.2)}
  \pgflineto{\pgfxy(0.9,0.1)}
  \pgflineto{\pgfxy(1,0)}
  \pgfclosepath
  \pgffill

  \color{structure}  
  \pgfputat{\pgfxy(0.25,0.7)}{\pgfbox[center,base]{#1}}
  \pgfputat{\pgfxy(0.25,-0.1)}{\pgfbox[center,top]{#2}}

  \color{black}
  \pgfputat{\pgfxy(-.1,0.25)}{\pgfbox[left,center]{\texttt{#3}}}%
}

\pgfdeclareverticalshading{heap1}{\the\paperwidth}%
  {color(0pt)=(black); color(1cm)=(structure!65!white)}
\pgfdeclareverticalshading{heap2}{\the\paperwidth}%
  {color(0pt)=(black); color(1cm)=(structure!55!white)}
\pgfdeclareverticalshading{heap3}{\the\paperwidth}%
  {color(0pt)=(black); color(1cm)=(structure!45!white)}
\pgfdeclareverticalshading{heap4}{\the\paperwidth}%
  {color(0pt)=(black); color(1cm)=(structure!35!white)}
\pgfdeclareverticalshading{heap5}{\the\paperwidth}%
  {color(0pt)=(black); color(1cm)=(structure!25!white)}
\pgfdeclareverticalshading{heap6}{\the\paperwidth}%
  {color(0pt)=(black); color(1cm)=(red!35!white)}

\newcommand{\heap}[5]{%
  \begin{pgfscope}
    \color{#4}
    \pgfheappath{\pgfxy(0,#1)}{\pgfxy(-#2,0)}{\pgfxy(#2,0)}
    \pgfclip
    \begin{pgfmagnify}{1}{#1}
      \pgfputat{\pgfpoint{-.5\paperwidth}{0pt}}{\pgfbox[left,base]{\pgfuseshading{heap#5}}}
    \end{pgfmagnify}
  \end{pgfscope}
  %\pgffill
  
  \color{#4}
  \pgfheappath{\pgfxy(0,#1)}{\pgfxy(-#2,0)}{\pgfxy(#2,0)}
  \pgfstroke

  \color{white}
  \pgfheaplabel{\pgfxy(0,#1)}{#3}%
}


\pgfdeclareradialshading{graphnode}
  {\pgfpoint{-3pt}{3.6pt}}%
  {color(0cm)=(beamerexample!15);
    color(2.63pt)=(beamerexample!75);
    color(5.26pt)=(beamerexample!70!black);
    color(7.6pt)=(beamerexample!50!black);
    color(8pt)=(beamerexample!10!averagebackgroundcolor)}

\newcommand{\graphnode}[2]{
  \pgfnodecircle{#1}[virtual]{#2}{8pt}
  \pgfputat{#2}{\pgfbox[center,center]{\pgfuseshading{graphnode}}}
}

%\setbeamerfont{author}{size=\Huge} 
%\setbeamerfont{date}{size=\small} 
%\setbeamerfont{exampletext}{series=\bfseries}
%\setbeamerfont{frametitle}{series=\bfseries,size=\large}
%\setbeamerfont{framesubtitle}{series=\mdseries,size=\large}
%\setbeamerfont{section}{series=\bfseries} 
\setbeamerfont{title}{series=\bfseries}

\setbeamercovered{invisible}

\makeatother

\usepackage{babel}
\begin{document}
\title{Calculus III}
\subtitle{MTH 331~~|~~Section A~~|~~Fall 2012}
\institute{{\Large\textbf{Johnson C. Smith }}}
\author{Dr.\,James Ashe}
\date{{}}

\makebeamertitle
%%suppress auto date
\title{Calc III: 9.1 Sequences and Infinites Series}

\author{}

\institute{}

\subtitle{}

\makebeamertitle
\begin{frame}{Sequences}

\begin{definition}

A \textbf{sequence }$\{a_{n}\}$ is an ordered list of numbers of
the form $\{a_{0},a_{1},a_{2},\dots,a_{n},\dots\}$.
\end{definition}

\begin{exampleblock}<2->{The Fibonacci Sequence}

$\{0,1,1,2,3,5,8,13,21,34,\dots\}$ where $a_{0}=0$, $a_{1}=1$,
and $a_{n}=a_{n-1}+a_{n-2}$ for $n>2$.
\end{exampleblock}
\begin{itemize}
\item <3->$a_{n}$\,is the $n^{th}$ term in the sequence; above the $4^{th}$
term is $a_{4}=2$.
\item <4->A sequence is generated by a \textbf{recurrence relation }if
further terms are defined as a function of preceding terms (like in
the Fibonacci sequence).
\end{itemize}
\end{frame}

\begin{frame}{}

A sequence may also be defined by an \textbf{explicit formula}. 
\begin{itemize}
\item <2->[]\textbf{ex) }$a_{n}=1+\sin(n\frac{\pi}{2})\Rightarrow$\uncover<3->{$\ensuremath{\{a_{n}\}_{n=1}^{\infty}}=\{2,1,0,1,2,1,0,\dots\}$}.
\item <4->There is an explicit formula for the Fibonacci sequence....
\item <5->[]...$a_{n}=\frac{1}{\sqrt{5}}\left((\varphi)^{n}-(1-\varphi)^{n}\right)$
where $\varphi=\frac{1+\sqrt{5}}{2}$, the \emph{Golden Ratio!} 
\item <6->More fun facts: ${\displaystyle \lim_{n\rightarrow\infty}\frac{a_{n}}{a_{n-1}}=\{\frac{a_{n}}{a_{n-1}}\}_{n=2}^{\infty}=\varphi}$
where $\{a_{n}\}$ is the Fibonacci sequence!
\item <7->We say that \textrm{if ${\displaystyle \lim_{n\rightarrow\infty}a_{n}=L}$
for a unique number $L$ than $a_{n}$ }\textrm{\textbf{converges}}\textrm{.
\uncover<8->{Otherwise we say that $a_{n}$ \textbf{diverges}}}
\item <9->[]\textbf{ex) }$\{1,2,3,\dots\}\rightarrow$\uncover<10->{diverges.}
\item <11->[]\textbf{ex) }$\{1,2,1,2,1\dots\}\rightarrow$\uncover<12->{diverges.}
\item <13->Fun question. Is $\{0,1,2,3,\dots\}>\{1,2,3,\dots\}?$
\end{itemize}
\end{frame}

\begin{frame}{Series}

\begin{definition}

An \textbf{infinite series }is the sum of an (infinite) sequence;
It has the form: ${\displaystyle \sum_{k=1}^{\infty}a_{n}=a_{1}+a_{2}+\dots+a_{n\dots}}=S_{n}$.
As before, if ${\displaystyle \lim_{n\rightarrow\infty}{\displaystyle \sum_{k=1}^{\infty}a_{n}={\displaystyle \lim_{n\rightarrow\infty}}S_{n}=L}}$
then the series \textbf{converges}; otherwise it \textbf{diverges}.
\end{definition}

\begin{itemize}
\item <2->[]\textbf{ex) }(see pg 52) Let $S_{n}={\displaystyle \sum_{k=1}^{n}}\frac{1}{2^{k}}$.
\item <3->[]\hspace{1cm}So then $S_{3}=$\uncover<4->{$\frac{1}{2^{1}}+\frac{1}{2^{2}}+\frac{1}{2^{3}}=\frac{7}{8}$.}
\item <5->[]\hspace{1cm}What is ${\displaystyle \lim_{n\rightarrow\infty}S_{n}}$?
\end{itemize}
\end{frame}

\begin{frame}{}

\textbf{HW 9.1: }21, 55, 60, 66, 68 

\vspace{1cm}

Next time, 9.2: Sequences
\end{frame}

\end{document}
